\clearpage

\section {Analysis procedure}
\label{sec:oscfit}

%--------------------------------------------------------------------------------------------------------------------
\subsection {Fit method}
\label{sec:oscfit_method}

The goal of this analysis is to measure the oscillation parameters
$\sin^2 \theta_{13}$, $\sin^2 \theta_{23}$, $\Delta m^2_{32}$ and
search for CP violation by setting limits on $\delta_{CP}$
for a given hypothesis of mass hierarchy, which is fixed to either normal or inverted.
All the systematics, as well as the oscillation parameters which are not measured
in that particular analysis and treated as nuisance, are marginalised.
The treatment of the oscillation parameters is shown in \cref{tab:osc_params_joint_priors}.
The parameters of interest will be denoted with the vector $\vec{\theta}$,
since more than 1 parameter at a time can be measured.
$\vec{\theta}$ can include the combination of all the oscillation parameters listed above.

%--------------------------------------------------------------------------------------------------------------------

The {\em VALOR} analysis uses a binned likelihood-ratio method where
the estimation of $\vec{\theta}$ is obtained
by comparing the observed and predicted Super-K $E_{reco}$ vs $\theta_l$ spectra
for 1-ring events.

Let $N$ be the number of reconstructed energy and lepton angle bins and
${\bf a}$ be a $1 \times N_{s}$ - dimensional array of systematic (nuisance) 
parameters that affect the Super-K reconstructed energy spectrum prediction (e.g.
parameters that control the beam, neutrino interaction simulations, 
detector response and oscillations).


The log-likelihood function used in the analysis is

\begin{equation} \label{eq:chisquare}
%\displaystyle
%\chi^{2} 
-2 \; \textrm{ln} \; \lambda (\vec{\theta},{\bf a}) = 2 \cdot \sum_{i=0}^{N-1} 
   \Big( 
     n^{obs}_{i} \cdot \textrm{ln} (n^{obs}_{i}/n^{exp}_{i}) +
     (n^{exp}_{i}-n^{obs}_{i}) 
   \Big) 
\end{equation}
where
$n^{obs}_{i}$ is the observed number of events in the $i^{th}$ bin and
$n^{exp}_{i} = n^{exp}_{i}(\vec{\theta}; {\bf a})$ 
%$n^{exp}_{i} = n^{exp}_{i}(\sin^{2}\overline{\theta}_{23},|\overline{\Delta m}^{2}_{32}|; {\bf a})$ 
is the corresponding expected number of events.


The log-likelihood function obtained by adding a gaussian penatly term can be written as
\begin{equation} \label{eq:chisquare_penalty}
%\displaystyle
%\chi^{2} 
-2 \; \textrm{ln} \; \lambda (\vec{\theta},{\bf a}) = 2 \cdot \sum_{i=0}^{N-1} 
   \Big( 
     n^{obs}_{i} \cdot \textrm{ln} (n^{obs}_{i}/n^{exp}_{i}) +
     (n^{exp}_{i}-n^{obs}_{i}) 
   \Big) +
   ({\bf a}-{\bf a_{0}})^{T} \cdot {\bf C}^{-1} \cdot ({\bf a}-{\bf a_{0}})
\end{equation}
where
${\bf a_{0}}$ is an $1 \times N_{s}$ - dimensional array with the default values of 
the systematic parameters, ${\bf a}^{T}$ is the transpose of ${\bf a}$,
${\bf C}$ is the systematic parameter covariance matrix of dimension $N_{s} \times N_{s}$,
$n^{obs}_{i}$ is the observed number of events in the $i^{th}$ bin, and
$n^{exp}_{i}(\vec{\theta}; {\bf a})$ 
%$n^{exp}_{i} = n^{exp}_{i}(\sin^{2}\overline{\theta}_{23},|\overline{\Delta m}^{2}_{32}|; {\bf a})$ 
is the corresponding expected number of events.

The total likelihood is given by the product of the likelihood terms of each selected sample,

\begin{equation} \label{eq:totlike}
\lambda(\vec{\theta}; {\bf a}) = \prod_{s=1}^{N_{samples}} \lambda_s (\vec{\theta}; {\bf a})
\end{equation}
where $N_{samples}$ is the number of selected samples used in the analysis, or
by the sum of $-2 \; \textrm{ln} \; \lambda (\vec{\theta},{\bf a})$ from all the Super-K samples

\begin{equation} \label{eq:totlike_2}
-2 \; \textrm{ln} \; \lambda (\vec{\theta},{\bf a}) = \sum_{s=1}^{N_{samples}} -2 \; \textrm{ln} \; \lambda_s (\vec{\theta},{\bf a})
\end{equation}

The advantage of the used binned likelihood ratio method is that,
in the large-sample limit, the quantity 
$-2 \: \textrm{ln} \: \lambda (\vec{\theta}; {\bf a})$
%$-2 \: \textrm{ln} \: \lambda (\sin^{2}\overline{\theta}_{23},|\overline{\Delta m}^{2}_{32}|; {\bf a})$
in \cref{eq:chisquare} has a $\chi^{2}$ distribution and it can therefore be used as a goodness-of-fit test.

Values of systematic parameters that give a negative predicted number of events in any reconstructed energy bin in any interaction 
mode are not allowed. If a systematic throw gives a negative number of events, that toy experiment is rejected.

%--------------------------------------------------------------------------------------------------------------------                                                                                    

\subsection {Marginalization}
\label{sec:marginalization}

In the previous analyses,
except \cite{OurRun56AntiNueBarAppearanceAnalysis3Flavour},
 a fully frequentist approach, in which all nuisance parameters were profiled, was used.
Now a hybrid Bayesian-frequentist method is used. 
The nuisance parameters are marginalised over while the parameters of interest are minimized with respect to the likelihood distribution.
The marginalization of the nuisance parameters give a so-called marginal likelihood that depends on only
on a vector of parameters of interest, $\vec{\theta}$.

The marginal likelihood is derived using the following equation:

\begin{equation}\label{eq:likemarg}
\lambda_{marg} (\vec{\theta}) = \int_{A} \lambda (\vec{\theta};{\bf a}) \pi{({\bf a})} d {\bf a} = \frac{1}{n} \sum_{i=1}^{n} 
\lambda (\vec{\theta}; {\bf a}_i) 
\end{equation}
where ${\bf a}$ is the set of nuisance parameters in the parameter space $A$, $\pi{({\bf a})}$ is the prior of the nuisance parameters.
$n$ is the number of toys MC produced for the marginalization and
${\bf a_i}$ is the single set of nuisance parameters extracted in the ``i-th" toy MC.
Many sets of ${\bf a}$ should be extracted in order to effectively sample the space of the nuisance parameters.
The number of extractions depends on the number of nuisance parameters as well as the precision that is required by the measurement.
In all the analyses shown in this note $n = 10k$ toys MC were used for the marginalisation of the systematic parameters.

We use the marginal log-likelihood, which will be denoted as $\chi^2_{marg}$ from now on,
since it follows a $\chi^2$ distribution when the Poisson statistic converges to the Gaussian:

\begin{equation} \label{eq:loglikemarg}
\chi^2_{marg} (\vec{\theta}) = - 2  \ln \lambda_{marg} (\vec{\theta})  = -2 \ln \left(\frac{1}{n} \sum_{i=1}^{n} 
\lambda (\vec{\theta}; {\bf a}_i)  \right) 
\end{equation}
It is important to note that this numerical integration is a weighted average of the likelihood evaluated at the $n$ throws, 
not a weighted average of the log-likelihood.
%When marginalization is used the penalty term
%$ ({\bf a}-{\bf a_{0}})^{T} \cdot {\bf C}^{-1} \cdot ({\bf a}-{\bf a_{0}}) $
%For this reason a modified version of \cref{eq:chisquare}, which doesn't include the penalty term, is used.
During the marginalization process, random values of the nuisance parameters are drawn according to the prior constraints which come from fits to near detector data and control samples at Super-K. Therefore, each throw is already weighted based on the prior information, and as a result the nuisance parameter values are not included in the calculation of the likelihood.
%
%When marginalization is used the penalty term becomes useless, 
%since the values of the nuisance parameters are already extracted using the prior distributions
%which replace the penalty term
%$ ({\bf a}-{\bf a_{0}})^{T} \cdot {\bf C}^{-1} \cdot ({\bf a}-{\bf a_{0}})$
%that also becomes constant, independent of the parameter of interest.
%
For this reason \cref{eq:chisquare}, which doesn't include the penalty term, is used.

Sometimes the free parameters, which are not marginalised, 
are constrained with a penalty term.
In this case a modified version of \cref{eq:loglikemarg} is used:

\begin{equation} \label{eq:chisqmarg_with_penalty}
\chi^2_{constr} (\vec{\theta}) = \chi^2_{marg}(\vec{\theta}) + (\vec{\theta} - \vec{\theta_0})^T V^{-1} (\vec{\theta} - \vec{\theta_0}) 
\end{equation}
where $\vec{\theta_0}$ is the vector of nominal values of the parameters of interest,
$\vec{\theta}$ is the vector of its fitted values 
and $V$ is the covariance matrix which relates all the free parameters.
If the constrained parameter $\theta_{constr}$ is not correlated with the other parameters of interest,
$V$ is reduced to the uncertainty $\Delta \theta^{constr}$,
and \cref{eq:chisqmarg_with_penalty} becomes

\begin{equation} \label{eq:chisqmarg_with_penalty_uncorr}
\chi^2_{constr} (\vec{\theta}) = \chi^2_{marg}(\vec{\theta}) + \frac{ (\theta^{constr} - \theta^{constr}_0)^2 }{(\Delta \theta^{constr})^2}
\end{equation}
This is the case of $\sin^2 \theta_{13}$ when it is measured, i.e. it is parameter of interest,
but at the same time it is constrained with the measurement of the reactor experiments.
When a penalty term is applied to a parameter of interest it will be specified.

The toys MC used for the marginalization of the nuisance parameters are performed exactly in the same way as 
described in \cref{sec:systematics:inputs}, taking into account the correlations, but without statistical fluctuations.
First $n = 10k$ distributions of $E_{reco}$ vs $\theta_l$ for $e$-like and $E_{reco}$ only for $\mu$-like samples
are produced using the extracted systematic parameters for a given point of the parameter of interest grid $\vec{\theta}$
for $1 \times 10^{20}$ POT so that they may be scaled appropriately when comparing to datasets of differing POT.
%and normalised to $1 \times 10^{20}$ POT, in order to be used for fits at any different number of POT. 
Finally they are stored in output files which can be used for whatever needed fit.
A fit on a single dataset is performed as following:
\begin{itemize}
\item all the $n$ toy MC distributions are read from the files,
\item each distribution is normalised to the number of POT of the data set,
\item for each toy distribution \cref{eq:chisquare} is calculated and $\lambda(\vec{\theta},{\bf a}) = e^{- (-2 \; \textrm{ln} \; \lambda (\vec{\theta},{\bf a})) / 2}$ is obtained,
\item the $\chi^2$ value is finally obtained, using \cref{eq:loglikemarg}, for a given set of parameter values of interest, $\vec{\theta}$.
	If $\vec{\theta}$ is constrained, \cref{eq:chisqmarg_with_penalty_uncorr} is used.
\end{itemize}

This procedure is followed for each point in the $\vec{\theta}$ grid, where the binning 
is so fine that it does not impact on the precision of the measurement.
Finally the best fit is found by looking for the $\vec{\theta}$ grid point 
with the minimum $\chi^2$ value, $\chi^2_{bf}$.
This point corresponds to the best-fit parameters $\vec{\theta}_{bf}$. 
The minimization can be also performed by using a discrete descent gradient algorithm \cite{DescentGradient}
implemented in the VALOR framework that allows the minimum $\chi^2$ 
to be found without the need of calculating the full $\chi^2$ distribution across all grid points.
This method is employed when a large number of fits is required, such as the calculation of critical values
for Feldman-Cousins contours, as it drastically reduces the amount of CPU time required.
%This will be very useful when many fits are done and a lot of CPU time is required.
A validation of the minimization algorithm will be shown in \cref{sec:validation} %are we showing this in the note or should we reference a previous one?
for the analyses where it is used.


%--------------------------------------------------------------------------------------------------------------------                                                                                    

\subsection{Best-fit spectra with marginalization}
\label{sec:spectra_bestfit_method}

When profiling, the best fit spectra are defined by the MC prediction that gives the largest likelihood values.

%If profiling is used the nuisance parameters are fixed to their best-fit values
%and the computation of the best-fit spectra is straightforward.

However, since the nuisance parameters are marginalised, there does not exist a single set 
of parameters
which corresponds to the best fit. Therefore a new method has been developed to 
represent the best fit spectra which is consistent with the oscillation results. 
%However, since the nuisance parameters are marginalised,
%in order to get the best-fit spectra consistent with the oscillation results,
%the nuisance parameters cannot be set to their best-fit values and
%a new method must be developed.

%At the end of the fit the nuisance parameters have an effective value which corresponds
%to the weighted average over all the toys MC where the weights are 
%the likelihood values of each single toy given by \cref{eq:chisquare}.
%This becomes clearer when considering \cref{eq:likemarg}.

The following procedure is used.
From each toy MC $i$, 
used to marginalise the nuisance parameters, 
a different spectrum,
with 
%$N^{j}(\vec{\theta}_{bf}; \bf{a})$ 
$N^{j}(\vec{\theta}_{bf})$ 
number of events in each bin $j$, 
is produced
as well as a likelihood $\lambda (\vec{\theta}_{bf},\bf{a_i}) $ value, 
calculated using \cref{eq:likemarg},
where 
$\vec{\theta}_{bf}$ is the best-fit value of the free parameters vector
and 
$\bf{a_i}$
is the vector of nuisance parameters corresponding to the toy $i$.
The best-fit number of events in the bin $j$ of the spectrum is calculated as

\begin{equation} \label{eq:bestfit_spectrum_marg}
N^j_{marg} = \frac{ \sum_{i=1}^{n} \lambda (\vec{\theta}_{bf}; {\bf a_i}) \cdot N^{j}(\vec{\theta}_{bf}; {\bf a_i}) }{ \sum_{i=1}^{n} \lambda (\vec{\theta}_{bf}; {\bf a_i})  }
\end{equation}
where $n$ is the number of toys MC.

One can easily find a relation between the profiled and marginalised best-fit spectra.
In the limit of infinite number of toys MC, 
i.e. the full nuisance parameters space is scanned,
the likelihood can be written as
\begin{equation} \label{eq:likemarg_bf}
\lambda_{marg} (\vec{\theta}_{bf}) = \frac{1}{n} \sum_{i=1,~i \neq bf}^{n} \left[ \lambda (\vec{\theta}_{bf}; \bf{a_i}) + \lambda_{prof}  \right]
\end{equation}
where 
$\lambda_{prof} =  \lambda (\vec{\theta}_{bf}; \bf{a_{bf}})$ 
is the profiled best-fit likelihood. 
If we denote as
$N^j_{prof} = N_i^{j}(\vec{\theta}_{bf}; \bf{a_{bf}})$  
the number of events in the bin $j$ obtained with profiling,
\cref{eq:bestfit_spectrum_marg} can be written as

\begin{equation} \label{eq:bestfit_spectrum}
N^j_{marg} = \frac{ \sum_{i=1,~i \neq bf}^{n} \lambda (\vec{\theta}_{bf}; \bf{a_i}) \cdot N^{j}(\vec{\theta}_{bf}; \bf{a_i}) + \lambda_{prof} \cdot N^{j}_{prof}  }{ \sum_{i=1}^{n} \lambda (\vec{\theta}_{bf}; \bf{a_i})  }
\end{equation}
%In the gaussian regime, 
%i.e. the likelihood has a linear dependence from the parameters 
%that are also distributed as gaussian,
%marginalization and profiling gives exactly the same result.
%If the term $\sum_{i=1,~i \neq bf}^{n} \lambda (\vec{\theta}_{bf}; \bf{a_i}) \cdot N^{j}(\vec{\theta}_{bf}; \bf{a_i})$
%is symmetric around $\mathcal{L}_{prof} \cdot N^{j}_{prof}$,
%the marginalised spectrum converges to the one obtained using the profiling and
%$N^j_{marg} = N^{j}_{prof}$.






%--------------------------------------------------------------------------------------------------------------------                                                                                    

\subsection{Confidence intervals}
\label{sec:cl}

After the fit is completed, 
the full $\chi^2$ distribution is then shifted with respect to $\chi^2 (\vec{\theta}_{bf})$:
\begin{equation} \label{eq:deltachisq}
\Delta \chi^2 (\vec{\theta}) = \chi^2(\vec{\theta}) - \chi^2(\vec{\theta}_{bf})
\end{equation}
Confidence intervals are set on $\Delta \chi^2 (\vec{\theta})$.
If 

\begin{equation} \label{eq:condition_cl}
\Delta \chi^2(\vec{\theta}) > \Delta \chi^2_{crit}
\end{equation}
where $\Delta \chi^2_{crit}$ 
is the critical value for a certain confidence level $X\%$(CL),
that value of $\vec{\theta}$ is excluded at the $X\%$ CL.
In the gaussian approximation constant values of $\Delta \chi^2_{crit}$ can be used.
It will be called ``constant $\Delta \chi^2$ method".

However in some cases the gaussian approximation is not reliable
and the ``Feldman-Cousins method" \cite{FeldmanCousins} must be used.
In this method a different critical value for each point of the $\vec{\theta}$ grid 
%$\Delta \chi^2 (\theta_{i..j})_{crit}$
$\Delta \chi^2 (\vec{\theta})_{crit}$
is calculated.
The procedure consists of the following steps:

\begin{enumerate}

\item many toys MC (e.g. 10k) are produced 
	assuming the oscillation hypothesis of the grid point $\vec{\theta}$
 	and taking into account both the statistical and systematic uncertainties; 

\item for each toy first $\chi^2(\vec{\theta}_{fix})$ is obtained
	by minimizing $\chi^2 (\vec{\theta})$
	where the oscillation parameters $\vec{\theta}$ are kept 
	fixed to the values corresponding to the grid point  
	and the nuisance parameters are marginalised.
	Then $\chi^2 (\vec{\theta})$ is minimized, 
	leaving the oscillation parameters $\vec{\theta}$ as free and
	$\chi^2 (\vec{\theta}_{bf})$ is obtained. 
	
\item the test statistic is calculated as

	\begin{equation}
	\Delta \chi^2 = \chi^2(\vec{\theta}_{fix}) - \chi^2 (\vec{\theta}_{bf})
	\end{equation}
	
%\item $\Delta \chi^2 (\vec{\theta})$ is calculated using 
%	\cref{eq:loglikemarg} or~\ref{eq:chisqmarg_with_penalty_uncorr}.

\item The steps 2 and 3 are repeated for each toy MC 
	and $f(\Delta \chi^2)$, the distribution of $\Delta \chi^2 (\vec{\theta})$,
	%analogue to what is shown in fig.~\ref{fig:sterile_pvalue},	
	is obtained.

\item Several critical values for $X\%$ CL (e.g. 68\%, 90\%...),
	%$\Delta \chi^2 (\vec{\theta})_{crit}^{X\%CL}$,
 	$\Delta \chi^2 (\vec{\theta})_{crit}$,
	can be defined as:
	
	\begin{equation} \label{eq:def_crit_val}
	\Delta \chi^2_{crit}: \int_{-\infty}^{\Delta \chi^2_{crit}} f(\Delta \chi^2) d( \Delta \chi^2) = X\%
	\end{equation}		
	
	where $f(\Delta \chi^2)$ is the p.d.f. of $\Delta \chi^2$ and 
	$d( \Delta \chi^2)$ is the differential of the integral.
	
\item This procedure is repeated for all the points of the $\vec{\theta}$ grid
	and several critical values are computed
	
\end{enumerate}

Once a critical value, $\Delta \chi^2 (\vec{\theta})_{crit}$,
is computed for each oscillation hypothesis in the $\vec{\theta}$ grid,
the confidence intervals are set by the condition of \cref{eq:condition_cl}. \\


Several analyses are performed and confidence intervals are computed for:
\begin{itemize}

\item 
1D $\delta_{CP}$ with reactor constraint: 
the parameter of interest is $\delta_{CP}$ and
all the other oscillation parameters are marginalised;
for $\sin^2 \theta_{13}$, that is marginalised,
the gaussian prior of reactor measurements shown in \cref{tab:osc_params_joint_priors}, is used. 

\item 
1D $\delta_{CP}$ T2K only: 
the parameter of interest is $\delta_{CP}$ and
all the other oscillation parameters are marginalised;
for $\sin^2 \theta_{13}$, that is marginalised, 
the uniform prior, shown in \cref{tab:osc_params_joint_priors}, is used. 

\item 
1D $\sin^2 \theta_{13}$ T2K only: 
the parameter of interest is $\sin^2 \theta_{13}$ and
all the other oscillation parameters are marginalised;
%for $\sin^2 \theta_{13}$, that is a free parameter,
%the gaussian prior of reactor measurements shown in \cref{tab:marg_binnning}, is used. 

\item 
1D $\Delta m^2_{32}$ ($\Delta m^2_{13}$) with reactor constraint: 
the parameter of interest is $\Delta m^2_{32}$ ($\Delta m^2_{13}$) and
all the other oscillation parameters are marginalised;
for $\sin^2 \theta_{13}$, that is marginalised,
the gaussian prior of reactor measurements shown in \cref{tab:osc_params_joint_priors}, is used. 

\item 
1D $\sin^2 \theta_{23}$ with reactor constraint: 
the parameter of interest is $\sin^2 \theta_{23}$ and
all the other oscillation parameters are marginalised;
for $\sin^2 \theta_{13}$, that is marginalised,
the gaussian prior of reactor measurements shown in \cref{tab:osc_params_joint_priors}, is used. 

%

\item 
2D $\sin^2 \theta_{13}$ vs $\delta_{CP}$ with reactor constraint: 
the parameters of interest are $\sin^2 \theta_{13}$ and $\delta_{CP}$ while
all the other oscillation parameters are marginalised;
for $\sin^2 \theta_{13}$, that is a free parameter, 
the gaussian prior of reactor measurements shown in \cref{tab:osc_params_joint_priors}, 
is used, through a penalty term as in \cref{eq:chisqmarg_with_penalty_uncorr}.

\item 
2D $\sin^2 \theta_{13}$ vs $\delta_{CP}$ T2K only: 
the parameters of interest are $\sin^2 \theta_{13}$ and $\delta_{CP}$ while
all the other oscillation parameters are marginalised;
%for $\sin^2 \theta_{13}$, that is a free parameter, 
%the uniform prior, shown in \cref{tab:marg_binnning}, is used. 

\item 
2D $\sin^2 \theta_{23}$ vs $\delta_{CP}$ T2K only: 
the parameters of interest are $\sin^2 \theta_{23}$ and $\delta_{CP}$ while
all the other oscillation parameters are marginalised;
%for $\sin^2 \theta_{13}$, that is a free parameter, 
%the uniform prior, shown in \cref{tab:marg_binnning}, is used. 

\item 
2D $\Delta m^2_{32}$ ($\Delta m^2_{13}$) vs $\delta_{CP}$ T2K only: 
the parameters of interest are $\Delta m^2_{32}$ and $\delta_{CP}$ while
all the other oscillation parameters are marginalised;
%for $\sin^2 \theta_{13}$, that is a free parameter, 
%the uniform prior, shown in \cref{tab:marg_binnning}, is used. 

\item 
2D $\sin^2 \theta_{23}$ vs $\Delta m^2_{32}$ ($\Delta m^2_{13}$) with reactor constraint: 
the parameters of interest are $\sin^2 \theta_{23}$ and $\Delta m^2_{32}$ ($\Delta m^2_{13}$) 
while all the other oscillation parameters are marginalised;
for $\sin^2 \theta_{13}$, that is marginalised,
the gaussian prior of reactor measurements shown in \cref{tab:osc_params_joint_priors}, is used. 

\end{itemize}

Each confidence interval is calculated for both normal and inverted hierarchy.

The 1D $\chi^2$ confidence intervals with ``constant $\Delta \chi^2$ method" are used to set the $\pm 1 \sigma$ ranges.

Each point of the $\vec{\theta}$ grid corresponds to the center of the histogram bins,
i.e. the $\chi^2$ values are calculated at each histogram bin center. 
The ranges and binnings used for the analyses listed above are shown in \cref{tab:marg_binnning}.
All the 1-dimensional analyses, where the precision on the best-fit point is important 
and must be quoted, are performed with a fine binning.
Since 2-dimensional analyses are performed only to set contours with ``constant $\Delta \chi^2$ method"
and a fit with a very high precision is not needed to find the best-fit point, 
coarser binning are used, without affecting the contour.
Furthermore a very fine 2-dimensional binning would require much more CPU time and storage space.
In the future $\bar{\nu}/\nu$ joint analysis releases when the "Feldman-Cousins method" will be needed 
also for 2-dimensional fits, a finer binning will be used.

\begin{table}[htb]
\centering
\begin{tabular}{| c | c | c | c | c |}
\hline
Parameter 					& Min 			& Max			& Step			& Bins	  	\\
\hline

& \multicolumn{4}{c|}{\bf 1-dimensional fit}  \\
\hline
$\delta_{CP}$ 					& $-\pi$			& $+\pi$			& $\pi / 200$		& 201		\\
$\sin^2 \theta_{13}$ 	(reactors)	 	& 0.015			& 0.036			& 0.021 / 81		& 201 		\\
$\sin^2 \theta_{13}$ 	(T2K only) 	& 0				& 0.113			& 0.113 / 200		& 201 		\\
$\sin^2 \theta_{23}$ 				& 0.3				& 0.7				& 0.4 / 200 		& 201		\\
$|\Delta m^2_{32}|~(\text{eV}^2 / \text{c}^4)$ 			& $2\times10^{-3}$	& $3\times10^{-3}$	& $10^{-3} / 200$	& 201 		\\
\hline

& \multicolumn{4}{c|}{\bf 2-dimensional fit}  \\
\hline
$\delta_{CP}$ 					& $-\pi$			& $+\pi$			& $\pi / 50$		& 51			\\
$\sin^2 \theta_{13}$ 	(reactors)	 	& 0.015			& 0.036			& 0.021 / 81		& 81	 		\\
$\sin^2 \theta_{13}$ 	(T2K only) 	& 0				& 0.113			& 0.113 / 81		& 81 			\\
$\sin^2 \theta_{23}$ 				& 0.3				& 0.7				& 0.4 / 81 			& 81			\\
$|\Delta m^2_{32}|~(\text{eV}^2 / \text{c}^4)$ 			& $2\times10^{-3}$	& $3\times10^{-3}$	& $10^{-3} / 21$	& 21 			\\
\hline

\end{tabular}
\caption{
Grid limits and spacing used for all the 1-dimensional and 2-dimensional marginalization fits and confidence regions.
The minimum and maximum values correspond to the center of the $\chi^2$ histogram first and last bin.
}
\label{tab:marg_binnning}
\end{table}





%--------------------------------------------------------------------------------------------------------------------                                                                                    

\subsection{Contours with profiled mass hierarchy}
\label{sec:cl_mh}

As already pointed out, the constant $\Delta \chi^2$ method gives the correct coverage 
when the gaussian approximation is valid and is very useful because it does not require a large amount of CPU time.
However, we know that in some cases the constant $\Delta \chi^2$ method fails to give the correct coverage
and can provide only approximate confidence intervals.
An example of this is the case when the constraint on a parameter is not precise enough to keep it away from a physical boundary. 
For cases such as these a toy MC study is required in order to estimate the correct critical values
(see \cref{sec:cl}).

Also, when considering discrete parameters, such as the mass hierarchy, that is distributed as binomial,
the $\Delta \chi^2$ is not expected to follow a 
$\chi^2$ distribution but more likely to follow a gaussian distribution as found in \cite{cl_masshierarchy}. This is an important 
distinction as the exclusion significance one can interpret from a constant $\Delta \chi^2$ contour is heavily dependant on the 
distribution the $\Delta \chi^2$ is expected to take. 

%\begin{equation}
%\sigma_{\Delta \chi^2} = 2 \sqrt{\Delta \chi^2}
%\label{eq:sigma_mh}
%\end{equation}
%which would not show the factor 2 in the gaussian approximation.

This means that when the constant $\Delta \chi^2$ method is used in order 
to set limits as a function of some oscillation parameters as well as the mass hierarchy,
the constant critical values we usually use
do not give the right coverage for the disfavoured mass hierarchy hypothesis contour.
%Indeed, since the $\sigma_{\Delta \chi^2}$ value that defines the significance is larger by a factor 2, 
%the gaussian approximation would erroneously produce 
%stronger limits, giving undercoverage.

For this reason the data fit 2-dimensional confidence intervals, 
which are calculated for both the hierarchy hypotheses, 
will be computed in two different ways:
\begin{itemize}
\item normal and inverted hierarchy contours are obtained independently
	with respect to their own minimum (``independent")
\item the global minimum between the normal and inverted hierarchy hypotheses
	is used to set both the contours (``global");
\end{itemize}
These different ways to show the confidence intervals correspond
to the two extremes, where the ``independent" and ``global" contours 
provide respectively the most and the least conservative results.

Since for the 1-dimensional fits the $\Delta \chi^2$ distribution are shown,
there is no need to provide results using both ``independent" and ``global" approaches.
However the reader should remember that the constant critical values 
shown in the figures hold only for the $\Delta \chi^2$ distribution corresponding to the best-fit mass hierarchy hypothesis.
Indeed the constant $\Delta \chi^2$ method does not give the proper coverage for discrete parameters
and consequently the contour of the disfavoured hierarchy has bad coverage.
%For the disfavoured hypothesis \cref{eq:sigma_mh}
%can be used independently for each point of the free parameter grid.




%--------------------------------------------------------------------------------------------------------------------                                                                                    

%\subsection{Goodness of fit}
%\label{sec:method_gof}
%
%The goodness-of-fit tests are done by minimising the expression in \cref{eq:loglikemarg} for $\delta_{CP}$
%as free parameter with reactor constraint.
%%In this analysis, the same coarse binning is used for these tests as in our neutrino-mode ??-disappearance analyses [4], [5] and [6]. 
%The choice of this binning followed a decision by the oscillation analysis group that p-values should be calculated
%using the rate and the spectra, and that a very coarse binning scheme should be
%used to ensure that there are sufficient numbers of events in each bin; 
%this scheme has 5 bins from 0-0.4, 0.4-0.7, 0.7-1.0, 1.0-2.0 and 2.0-30.0 GeV. 
%%The Asimov spectrum using this binning can be seen in figure 23.
%1000 toy experiments are generated for each mass hierarchy, and both systematic variations and statistical
%fluctuations are applied to them. All \NSyst{} systematic parameters are included in the fits of those datasets.
%
%
%%\vspace{3mm}
%%{\huge{
%%TODO: DESCRIBE HOW THE P-VALUE IS CALCULATED ???
%%}}
%%\vspace{3mm}


\newpage











